\section{Discussão dos Achados}
\label{sec:discussao_resultados}

Os testes verificam comportamentos especificados do código; o questionário
registra percepção e dificuldades declaradas. A discussão relaciona esses
resultados com a literatura e com as ações descritas nas tarefas. As limitações da
avaliação estão reunidas na Seção~\ref{sec:ameacas_validade}.

\subsection{Diagnósticos, ação de correção e iniciantes}

Os oito usuários concordaram com os itens sobre localização, causa provável e
próximo passo, enquanto seis dos oito concordaram que corrigiram sem explicação
adicional. Essa diferença distingue perceber uma mensagem como clara e conseguir
agir de forma independente. Na perspectiva de Marceau, Fisler e Krishnamurthi
(2011a), apresentada no Capítulo~\ref{cap:fundamentacao}, localizar, compreender e
formular uma correção são aspectos que não se confirmam apenas por uma avaliação
favorável do texto. Os relatos de T2 e T3 contêm exemplos compatíveis com correção,
mas não documentam integralmente essas etapas para cada pessoa.

Denny, Luxton-Reilly e Carpenter (2014) e Becker (2016) relataram resultados
diferentes ao avaliar melhorias em mensagens, em contextos e medidas próprios.
O estudo de Becker examinou erros de estudantes de Java, e não apenas preferência
pelas explicações. O Luminthor, por sua vez, não coletou logs equivalentes de
sucessivas tentativas nem contou com grupo comparador. Os achados de mensagens
mais claras são, portanto, indícios de aceitação a investigar por observação,
sem demonstrar a mesma redução de erros ou superioridade sobre outro ambiente.

O sinal menos favorável dos iniciantes, destacado na
Seção~\ref{sec:usabilidade_resultados}, deve orientar a próxima avaliação.
A confirmação do cruzamento original permanece pendente, mas o relato registrado
já recomenda examinar controles, termos e necessidade de explicação em vez de
apoiar a decisão apenas nas medianas gerais. Além disso, a heurística de nomes
curtos documentada na Seção~\ref{sec:recuperacao_sugestoes} mostra que uma sugestão
determinística pode ser pouco útil. Isso qualifica a opção por sugerir correções
frente à cautela de Marceau, Fisler e Krishnamurthi (2011b): a sugestão é uma
possibilidade a avaliar, não uma reconstrução da intenção do aprendiz.

\subsection{Depurador, máquina nocional e evidência das tarefas}

U07 recebeu oito respostas favoráveis, sete delas com nota 5. O motor também
produz estados e rastro verificáveis, como no exemplo da
Seção~\ref{sec:exemplo_pipeline}. Entretanto, os relatos de T1 não distinguem em
um mesmo registro a saída e a leitura das variáveis, e nenhum relato de T3
explicita o estado anterior à falha. A preferência pelo depurador e sua existência
técnica não confirmam que os participantes tenham interpretado os estados solicitados.

Sorva, Karavirta e Malmi (2013) distinguem formas de envolvimento com a visualização
e apontam limites nas conclusões sobre seus efeitos. À luz dessa discussão e da
máquina nocional apresentada na Seção~\ref{sec:visualizacao_programas}, interessa
observar o que o iniciante faz com a representação: explicar um valor, prever o
próximo estado e relacionar uma mudança ao comando correspondente. Sessões observadas em trabalhos futuros podem solicitar essas ações, para
registrar a interpretação dos estados durante as tarefas.

\subsection{Interação móvel, carga cognitiva e comparação com ferramentas}

U08 também recebeu oito respostas favoráveis, sete com nota 5, mas um dos quatro
temas dos comentários trata da digitação em tela pequena. A avaliação global da
interface pode coexistir com dificuldade em uma ação específica. Cinco dos oito
usuários consultaram o formulário no mesmo celular, condição que pode exigir
alternância entre instruções e editor; esse registro não caracteriza dependência
exclusiva do aparelho para estudar.

A teoria de carga cognitiva discutida no Capítulo~\ref{cap:fundamentacao}, com
Sweller (1988) e Sweller, van Merriënboer e Paas (1998), oferece uma interpretação
para investigar navegação, nomenclatura e divisão da atenção. O relato sobre
Trace e as dificuldades de edição justificam observar essas ações. Contudo, não
foi aplicada medida de carga cognitiva: não se conclui que abas, atalhos ou
explicações tenham reduzido carga extrínseca. Expor muitos painéis também pode
exigir orientação inicial; essa possibilidade deve ser examinada com iniciantes.

O quadro comparativo da Seção~\ref{sec:comparacao_ferramentas} documenta recursos
próximos em outras ferramentas, incluindo execução e depuração em Portugol Mobile
e WebStudio, além de visualização em Python Tutor e Thonny. Os achados do Luminthor
não resultam de uma comparação experimental com esses ambientes. Seu resultado
tecnológico é a integração específica de linguagem, motor local, editor, inspeção
e trilha, com os compromissos descritos na Seção~\ref{sec:decisoes_tecnologicas}.
As classificações documentais ainda não confirmadas na comparação não sustentam
exclusividade ou superioridade do protótipo.

\begin{landscape}
\section{Síntese do Alcance das Evidências}
\label{sec:sintese_evidencias}

O Quadro~\ref{quad:alcance_evidencias} separa demonstração técnica, relato e
verificação pendente para as afirmações centrais. ``Demonstrado'' se refere ao
escopo do código e dos casos registrados. As páginas indicam onde a evidência ou sua pendência é apresentada.

\begin{quadro}[htb]
\centering
\legenda[quad:alcance_evidencias]{Demonstrado, relatado e não verificado no Luminthor}
\footnotesize
\renewcommand{\arraystretch}{1.12}
\begin{tabular}{|P{2.7cm}|P{6.2cm}|P{5.7cm}|P{7.4cm}|}
\hline
Afirmação & Demonstrado no escopo técnico & Relatado na avaliação & Não verificado \\
\hline\hline
Processamento local e acesso Android & Motor e chamadas FFI nos testes atuais, Seção~\ref{sec:verificacao_atual}, p.~\pageref{sec:verificacao_atual}. & Oito de 11 consentidos declararam abertura, Tabela~\ref{tab:perfil_participantes}, p.~\pageref{tab:perfil_participantes}. & Ensaio controlado em modo avião e compatibilidade em aparelhos, não realizados. \\
\hline
Diagnósticos em português & Programas e mensagens principais reproduzidos, Quadros~\ref{quad:regras} e~\ref{quad:diagnosticos_execucao}, p.~\pageref{quad:regras}. & Oito de oito favoráveis em U02, U03 e U05; seis em U06, Tabela~\ref{tab:usabilidade_participantes}, p.~\pageref{tab:usabilidade_participantes}. & Correção autônoma de cada tarefa e causa do relato divergente T3; Seção~\ref{sec:reproducao_t3}, p.~\pageref{sec:reproducao_t3}. \\
\hline
Rastro e inspeção de variáveis & Cinco eventos do exemplo, valores e saída, Quadro~\ref{quad:rastro_exemplo}, p.~\pageref{quad:rastro_exemplo}; limites de escopo na Seção~\ref{sec:semantica_lumi_detalhada}. & Oito de oito favoráveis em U07; registros parciais de T1/T3, Tabelas~\ref{tab:usabilidade_participantes} e~\ref{tab:respostas_tarefas}, p.~\pageref{tab:respostas_tarefas}. & Interpretação independente dos estados e aprendizagem, Seção~\ref{sec:discussao_resultados}, p.~\pageref{sec:discussao_resultados}. \\
\hline
Edição, persistência e arquivos & Casos de widgets e regras; limitação de falha do salvamento descrita no Desenvolvimento, Seção~\ref{sec:categorias_testes}, p.~\pageref{sec:categorias_testes}. & Oito de oito favoráveis em U08, com dificuldade de digitação entre os temas, Quadro~\ref{quad:temas_comentarios}, p.~\pageref{quad:temas_comentarios}. & Persistência após reabertura, importação, exportação e teclado em ensaio controlado em aparelhos, não realizado. \\
\hline
Trilha com 18 aulas & Conteúdo, sequência e testes das regras de progressão, Quadro~\ref{quad:aulas_lumi}, p.~\pageref{quad:aulas_lumi}. & Sem resultado específico de eficácia da trilha registrado no levantamento. & Aprendizagem, transferência e conclusão independente dos desafios, p.~\pageref{sec:ameacas_validade}. \\
\hline
Adequação a iniciantes & Artefato implementado e documentado. & Dois usuários com menos de três meses, autores das respostas neutras em U01, U04 e U06, Seção~\ref{sec:usabilidade_resultados}, p.~\pageref{sec:usabilidade_resultados}. & Observação com iniciantes e benefício educacional. \\
\hline
Desempenho móvel & Contagens e cobertura de linhas são verificações de código, p.~\pageref{tab:verificacao_complementar}. & As respostas de usabilidade não medem tempos ou memória. & Tempos, memória, tamanho e abertura do APK, Tabelas~\ref{tab:medicoes_motor} e~\ref{tab:medicoes_apk}, p.~\pageref{tab:medicoes_motor}. \\
\hline
\end{tabular}
\vspace{0.2cm}
\fonte{O autor (2026).}
\end{quadro}
\FloatBarrier
\end{landscape}

\section{Contribuição e Continuidade da Avaliação}
\label{sec:contribuicao_resultados}

A contribuição tecnológica reúne o artefato integrado, suas decisões de projeto,
a especificação e os casos de teste. As Seções anteriores identificam as evidências
disponíveis e as verificações necessárias para dar continuidade à avaliação.

Entre os oito usuários, houve de seis a oito respostas favoráveis por item.
Os registros indicam prioridades concretas: orientação inicial, nomenclatura do
depurador, edição em telas pequenas e verificação do caso T3. A continuidade deve acompanhar essas ações e registrar as evidências que faltam.
